\chapter{Number Theory}
Well-Ordering Principle: Every non-empty subset of non-negative integers set has a minimum element.
\section{Definition}
\begin{definition}
    Let $a, b \in \mathbb{Z}$ be integers. \\
    We said to be $b$ divides $a$ if: $a = bx$ for some $x \in \mathbb{Z}$. Denote $b \mid a$.
\end{definition}

\section{The Division Algorithm}
\begin{theorem}
    Let $a, b \in \mathbb{Z}$ be integers with $b > 0$. 
    Then, there exists unique integers $q, r \in \mathbb{Z}$ such that
    \begin{align*}
        a = b \cdot q + r \quad (0 \leq r < b)
    \end{align*}
\end{theorem}
\begin{proof}
    Let $S \defequal \{a - bx \mid x \in \mathbb{Z},\ a - bx \geq 0\}$.
    Then, $S$ is non-empty because $a - b \cdot (-|a|) = a + |a| b \geq a + |a| \geq 0$. \\
    Thus, Well-Ordering Principle gives that there exists a minimum element of $S$. Put $r = \min S$. \\
    (That is, for any $k \in S$, $r \leq k$.)
    Since $r \in S$, there exists $q \in \mathbb{Z}$ such that $ r = a - bq \geq 0$. \\
    To show $r < b$, we will use Contradiction. Suppose that $r \geq b$. Then,
    \begin{align*}
        a - b \cdot (q + 1) = a - bq - b = r - b \geq 0
    \end{align*}
    Thus, $a - b (q + 1) \in S$, but this is contradiction with $r = a - bq$ is a minimum element of $S$. \\
    To show the Uniqueness, there exist $q, q', r, r' \in \mathbb{Z}$ such that
    \begin{align*}
        a = bq + r = bq' + r' \quad(0 \leq r, r' < b)
    \end{align*}
    Then, $b|q - q'| = |r - r'|$, thus must be $q = q'$ and $r = r'$.
\end{proof}
\begin{corollary}
    Let $a, b \in \mathbb{Z}$ be integers with $b \neq 0$. 
    Then, there exists unique integers $q, r \in \mathbb{Z}$ such that
    \begin{align*}
        a = b \cdot q + r \quad (0 \leq r < |b|)
    \end{align*}
\end{corollary}

\newpage
\section{Bézout's Identity}
\begin{definition}
    Let $a,b \in \mathbb{Z}$ be integers. The integer $d$ is called \textit{The greatset common divisor} of $a, b$ if:
    \begin{enumerate}
        \item $d \mid a$ and $d \mid b$.
        \item If $c \mid a$ and $c \mid b$, then $c \leq d$.
    \end{enumerate}
    Denote $d \defequal \gcd(a,b)$.
\end{definition}

\begin{theorem}
    Let $a,b \in \mathbb{Z}$ be integers. Then, there exist integers $x, y \in \mathbb{Z}$ such that
    \begin{align*}
        \gcd(a,b) = ax + by
    \end{align*}
\end{theorem}
\begin{proof}
    Define a set:
    \begin{align*}
        S \defequal \{a \cdot u + b \cdot v \mid u,v \in \mathbb{Z},\ a \cdot u + b \cdot v > 0\}
    \end{align*}
    Clearly, $S$ is non-emptyset. Thus, $S$ has a minimum element $d = \min S$ by Well-Ordering Principle. \\
    Since $d \in S$, $d = ax + by$ for some $x,y \in \mathbb{Z}$. By the Division Algorithm, there exist $q, r \in \mathbb{Z}$ such that
    \begin{align*}
        a = d \cdot q + r \quad(0 \leq r < d)
    \end{align*}
    That is,
    \begin{align*}
        r = a - d \cdot q = a - q \cdot (ax + by) = a \cdot (1 - qx) + b \cdot (-qy)
    \end{align*}
    If $r \neq 0$, then $r \in S$ and $r < d$, this is contradiction with $d = \min S$. Thus, $r = 0$. \\
    Similarly, $b = d \cdot q'$ for some $q' \in \mathbb{Z}$. In summary, $d$ divides $a$ and $b$. \\
    Now, let $c$ be an arbitrary positive common divisor of $a$ and $b$.
    Then, $c$ divides $ax + by = d$, thus $c \leq d$. \\
    Consequently, $d = \gcd(a,b)$.
\end{proof}
\begin{corollary}
    Let $a,b$ be integers. For $d = \gcd(a,b)$,
    \begin{align*}
        \{a \cdot x + b \cdot y \mid x, y \in \mathbb{Z}\} = \{d \cdot k \mid k \in \mathbb{Z}\}
    \end{align*}
\end{corollary}

\newpage
\subsection{Extended Euclidean Algorithm}
\begin{lemma}
    Let $a,b$ be integers. \\
    If $a = b \cdot q + r$ with $0 \leq r < b$,
    then $\gcd(a,b) = \gcd(b,r)$.
\end{lemma}
\begin{proof}
    Put $d = \gcd(a,b)$. Since $d \mid a$ and $d \mid b$, $d$ divides $r = a - bq$.
    That is, $d$ is common divisor of $b$ and $r$. \\
    Meanwhile, let $c$ be a divisor of $b$ and $r$. Then, $c$ divides $a = bq + r$, thus $c \leq d$.
    This implies $d = \gcd(b,r)$.
\end{proof}
Consider the given integers $a,b$. By the Division Algorithm, we can write:
\begin{align*}
    a &= b \cdot q_0 + r_0 &&(0 \leq r_0 < b)\\
    b &= r_0 \cdot q_1 + r_1 &&(0 \leq r_1 < r_0)\\
    r_0 &= r_1 \cdot q_2 + r_2 &&(0 \leq r_2 < r_1)\\
    &\ \vdots \\
    r_{k-2}&= r_{k-1} \cdot q_{k} + r_{k} &&(0 \leq r_k < r_{k-1})\\
    &\ \vdots \\
    r_{n-2} &= r_{n-1} \cdot q_n + r_n &&(0 \leq r_n < r_{n-1})\\
    r_{n-1} &= r_{n} \cdot q_{n+1} &&(r_{n+1} = 0)
\end{align*}
Since the Non-negative numbers Sequnece $\{r_k\}$ strictly decrease, $r_{n+1}$ must be zero for some $n \in \mathbb{N}$. \\
Put $r_{-2} = a$ and $r_{-1} = b$. Since above Lemma,
\begin{align*}
    \gcd(a,b) = \gcd(b, r_0) = \gcd(r_0, r_1) = \cdots = \gcd(r_{n-1}, r_n) = \gcd(r_n, 0) = r_n
\end{align*}
This implies also for any $-2 \leq k \leq n$, $r_k$ is multiple of $r_n$. \\
% By Bézout’s Identity, there exists $x, y \in \mathbb{Z}$ such that $r_n = \gcd(a,b) = a \cdot x + b \cdot y$. \\
% Meanwhile, for $k = n$, $r_{n-1} = r_n \cdot q_{n+1}$ is multiple of $r_n$.
% Put $r_{-2} = a$ and $r_{-1} = b$. Suppose that for $0 \leq k \leq n$, $r_{k-1}$ and $r_k$ both are multiple of $r_{n}$. Then, the formular
% \begin{align*}
%     r_{k-2} = r_{k-1} \cdot q_n + r_k
% \end{align*}
% implies $r_{k-2}$ is multiple of $r_n$. And, for $k = n$, $r_{n-1} = r_n \cdot q_{n+1}$ is multiple of $r_n$. \\
% Now, the Induction implies for all $0 \leq k \leq n$, $r_k$ is multiple of $r_n$. \\
By Corollary of Bézout’s Identity, for any $-2 \leq k \leq n$, there exists $x_k, y_k \in \mathbb{Z}$ such that
\begin{align*}
    r_k = a \cdot x_k + b \cdot y_k \tag{$*$}
\end{align*}
Since $r_{-2} = a$ and $r_{-1} = b$,
\begin{align*}
    &x_{-2} = 1,\ y_{-2} = 0 \\
    &x_{-1} = 0,\ y_{-1} = 1
\end{align*}
And, substituting $(*)$ to $r_{k-2} - r_{k-1} \cdot q_k = r_k$:
\begin{align*}
    a \cdot x_{k} + b \cdot y_{k}
    = a \cdot x_{k-2} + b \cdot y_{k-2} - (a \cdot x_{k-1} + b \cdot y_{k-1}) \cdot q_{k}
    = a \cdot (x_{k-2} - x_{k-1} \cdot q_k) + b \cdot (y_{k-2} - y_{k-1} \cdot q_k)
\end{align*}
Now, we obtain the equation: for all $0 \leq k \leq n$,
\begin{align*}
    &x_k = x_{k-2} - x_{k-1} \cdot q_k \\
    &y_k = y_{k-2} - y_{k-1} \cdot q_k
\end{align*}
Consequently, to find $x,y \in \mathbb{Z}$ such that $\gcd(a,b) = r_n = a \cdot x + b \cdot y = a \cdot x_n + b \cdot y_n$, \\
we simply calculate $x_n, y_n$ inductively.

% \begin{proof}
    % By the Division Algorithm, there exist $q_0, r_0$ such that $a = b \cdot q_0 + r_0$ with $0 \leq r_0 < b$. \\
    % Using again Division Algorithm for $b$ and $r_0$, there exist $q_1, r_1$ such that $b = r_0 \cdot q_1 + r_1$ with $0 \leq r_1 < r_0$. 
    % Inductively,
    % \begin{align*}
    %     a &= b \cdot q_0 + r_0 \\
    %     b &= r_0 \cdot q_1 + r_1 \\
    %     r_0 &= r_1 \cdot q_2 + r_2 \\
    %     & \vdots \\
    %     r_{k-2}&= r_{k-1} \cdot q_{k} + r_{k} \\
    %     & \vdots \\
    %     r_{n-2} &= r_{n-1} \cdot q_n + r_n \\
    %     r_{n-1} &= r_{n} \cdot q_{n+1}
    % \end{align*}
    % Since the Non-negative numbers Sequnece $\{r_k\}$ strictly decrease, $r_{n+1}$ must be zero for some $n \in \mathbb{N}$. \\
    % And, since above Lemma,
    % \begin{align*}
    %     \gcd(a,b) = \gcd(b, r_0) = \gcd(r_0, r_1) = \cdots = \gcd(r_{n-1}, r_n) = \gcd(r_n, 0) = r_n
    % \end{align*}
    % Now, rewrite the above formulas:
    % \begin{align*}
    %     r_n &= r_{n-2} - r_{n-1} \cdot q_n\\
    %     &= r_{n-2} - (r_{n-3} - r_{n-2} \cdot q_{n-1}) = r_{n-2} \cdot (1 + q_{n-1}) - r_{n-3} \\
    %     &= (r_{n-4} - r_{n-3} \cdot q_{n-2}) \cdot (1 + q_{n-1}) - r_{n-3}
    %     = r_{n-4} \cdot (1 + q_{n-1}) - r_{n-3} \cdot (1 + q_{n-2} + q_{n-1} \cdot q_{n-2}) \\
    %     &= r_{n-4} \cdot (1 + q_{n-1}) - (r_{n-5} - r_{n-4} \cdot q_{n-3}) \cdot (1 + q_{n-2} + q_{n-1} \cdot q_{n-2}) \\
    %     &= r_{n-4} \cdot (1 + q_{n-1} + q_{n-3} + q_{n-2} \cdot q_{n-3} + q_{n-1} \cdot q_{n-2} \cdot q_{n-3}) - r_{n-5} \cdot (1 + q_{n-2} + q_{n-1} \cdot q_{n-2}) \\
    %     &= (r_{n-6} - r_{n-5} \cdot q_{n-4}) \cdot (1 + q_{n-1} + q_{n-3} + q_{n-2} \cdot q_{n-3} + q_{n-1} \cdot q_{n-2} \cdot q_{n-3}) - r_{n-5} \cdot (1 + q_{n-2} + q_{n-1} \cdot q_{n-2}) \\
    %     &= r_{n-6} \cdot (1 + q_{n-1} + q_{n-3} + q_{n-2} \cdot q_{n-3} + q_{n-1} \cdot q_{n-2} \cdot q_{n-3})
    %     - r_{n-5} \cdot (1 + q_{n-2} + q_{n-1} \cdot q_{n-2} )
    % \end{align*}
    
    % \begin{align*}
    %     a - b \cdot q_0 &= r_0 \\
    %     b - r_0 \cdot q_1 &= r_1 \\
    %     r_0 - r_1 \cdot q_2 &= r_2 \\
    %     & \vdots \\
    %     r_k - r_{k+1} \cdot q_{k+2} &= r_{k+2} \\
    %     & \vdots \\
    %     r_{n-4} - r_{n-3} \cdot q_{n-2} &= r_{n-2} \\
    %     r_{n-3} - r_{n-2} \cdot q_{n-1} &= r_{n-1} \\
    %     r_{n-2} - r_{n-1} \cdot q_n &= r_n
    % \end{align*}
    % Substituting gives: for $n \geq 2$,
    % \begin{align*}
    %     r_0 &= a - b \cdot q_0 \\
    %     r_1 &= b - r_0 \cdot q_1 = b - (a - b \cdot q_0) \cdot q_1 = a \cdot (-q_1) + b(1 + q_0 \cdot q_1) \\
    %     r_2 &= r_0 - r_1 \cdot q_{2} = a - b \cdot q_0 - (a \cdot (-q_1) + b(1 + q_0 \cdot q_1)) \cdot q_2
    %     = a \cdot(1 + q_1 \cdot q_2) - b \cdot (q_0 + )
    % \end{align*}
% \end{proof}